\documentclass[sigconf,anonymous,review]{acmart}

\input{macros}

\setcopyright{none}
\settopmatter{printacmref=false}
\renewcommand\footnotetextcopyrightpermission[1]{}
\acmConference{}{}{}
\acmYear{}
\acmDOI{}
\acmISBN{}

\begin{document}

\title{Email Recovery Without Plaintext: Zero-Knowledge Email as a Recovery Factor for Smart Accounts}

\author{Anonymous Author(s)}
\affiliation{%
  \institution{Anonymous Institution}
  \country{Anonymous Country}}

\begin{abstract}
A smart account can use control of an email inbox as one recovery factor
without anyone storing a plaintext email address or secret. The account
commits to a salted hash of the address. To recover, the user sends one email
from their own inbox carrying the recovery statement, proves in the browser
that the provider DKIM-signed it, and an on-chain verifier accepts the proof
as one guardian approval inside an existing delayed recovery process. Email is
never sufficient alone. We give the design on Perch's recovery for Soroban
smart accounts, state its soundness, privacy, replay-resistance and
one-factor-safety properties, and mark each claim as proven, tested,
measured, assumed or unproven. Implementation and evaluation results are
pending.
\end{abstract}

\keywords{account recovery, DKIM, zero-knowledge proofs, smart accounts, Soroban}

\maketitle

\section{Introduction}\label{sec:intro}

A smart account that authenticates with a passkey is lost when every device
holding the passkey is lost. Account-abstraction wallets answer this with
recovery: guardians who approve a key change, a delay during which the owner
can object, and secondary credentials~\cite{buterin-social-recovery,argent-spec,erc4337}.
An email inbox is the secondary credential most people already have and keep
for years, and large services have long used it to recover
accounts~\cite{bonneau-recovery}. Used naively, it puts the address in the
operator's database and makes whoever runs the email check a party that can
recover any account: a server that mails a one-time code, or a recovery
server that holds signing keys for its users~\cite{sep30}, is trusted with
every account it serves. Checking the DKIM signature~\cite{rfc6376} on chain
instead removes the server but publishes the signed \texttt{From} header.

This paper shows that control of an inbox can be one recovery factor while
the chain holds only a salted hash of the address,
$\Commit$, where $\code$ is an account code that only the user holds, and
the operator's servers see the address only inside an attested enclave. To
recover, the user sends one email from the inbox to a fixed address. Its
Subject carries a blinded copy of the digest of the pending recovery
statement, so the email's content names no account or attempt. The user's
browser proves in zero knowledge that the provider DKIM-signed the email
under a registered key, that its sender opens the commitment, and that its
Subject opens to the statement's digest. A per-account email-guardian
contract verifies the proof on Soroban and counts it as one guardian
approval in Perch's recovery~\cite{perch-recovery-spec}, used unchanged.
Perch's quorum and delay still apply: email approvals never complete a
recovery before the delay has passed, and they authorize one alone only
under a lost-key profile with at least seven days of delay and a second
alert channel, where the owner's cancellation is the defence.

Five problems shape the design. \emph{Binding}: an email must approve exactly
one recovery statement (one account, new key, attempt and deadline), and a
copy must be worth no more than the original. \emph{Parsing soundness}: a
circuit must read a header the way the provider signed it; public audits of
open-source email circuits have repeatedly found sender spoofing and
header-confusion
flaws~\cite{audit-zksecurity-2024,audit-yacademy,audit-veridise-2024,audit-consensys-2024,audit-matterlabs-2024,audit-zellic-2024}.
\emph{Key trust}: whoever says which RSA key signs for \texttt{gmail.com}
must not be able to plant a key and recover accounts with it. \emph{Privacy
from the operator}: the operator's servers and staff must not read the
address at enrollment, at recovery, or when sending alerts. \emph{The
human}: a user can be talked into sending the email and reading out the
code, and under email plus delay only an alert that reaches the owner in
time stops a wrongful recovery.

\paragraph{Contributions.}
\begin{itemize}
  \item A design for email as one guardian inside an existing recovery
    protocol, unchanged (\S\ref{sec:design}). What is new against zkEmail's
    recovery stack (\S\ref{sec:related}): the Subject binds the protocol's
    own statement digest, blinded, rather than a readable command; the proof
    is made in the user's browser, so no relayer receives the email or holds
    the account code; the key registry is immutable, fed by DNS attestations,
    with public timelocks, and the per-account pin is salted, expiring and
    fixed with the guardian; and alerts come from the protocol's own events
    through a sealed vault, with a liveness probe and a wallet that reads
    open attempts itself\claim{guardian-on-perch}.
  \item Four properties stated over a Lean model of the guardian, the
    circuit's relation and Perch's controller (\S\ref{sec:properties}).
    \emph{Proven} over the model: soundness and replay resistance, each the
    property or a hash collision on inputs the theorem computes, with
    knowledge extraction per accepted proof; one-factor safety over every
    trace; privacy as a simulation theorem in an ideal model of the commitment
    and the
    proofs\claims{email-factor-soundness,replay-resistance,one-factor-safety,address-privacy}.
    \emph{Assumed}: the cryptography behind the ideal model and the named
    breaks, and that providers authenticate senders. \emph{Tested}: the
    model's agreement with the contracts, and the circuit's parsing.
  \item A Noir circuit that refuses each publicly documented soundness issue
    class of open-source email circuits at a named assertion, and whose
    parsing agrees with reference implementations on real mail and fuzzed
    headers\claims{circuit-issue-classes,parser-real-mail,parser-fuzzing}.
  \item An evaluation on the test host, a local network running mainnet's
    settings, and testnet (\S\ref{sec:eval}): a 73-row attack suite that
    includes phishing, deleted alerts, a lapsed push, a watcher outage and
    repeated attempts; each environment's fees attributed to its settings;
    archival and its keeper; proving time on a desktop; phones not measured.
  \item A claim ledger giving every claim a status and its evidence, checked
    by the build (Appendix~\ref{app:claims}).
\end{itemize}

\paragraph{Claim tags.} A tag ends each statement about the system:
\claimtag{proven} for a Lean theorem about the model, \claimtag{tested} for
named tests, \claimtag{measured} for a recorded measurement,
\claimtag{assumed} for a stated assumption, \claimtag{unproven} for a
statement without evidence. Appendix~\ref{app:claims} lists each claim with
its evidence; the build fails if a tag names no ledger entry, if evidence is
missing, or if a measured figure does not occur in its measurement files.

\section{Background}\label{sec:background}

\subsection{DKIM}
DomainKeys Identified Mail~\cite{rfc6376} lets the organisation that owns a
signing domain take responsibility for a message. The signer adds a
\texttt{DKIM-Signature} header field naming the signing domain
(\texttt{d=}), a selector (\texttt{s=}) under which the public key is
published in DNS at \texttt{s.\_domainkey.d}, the list of signed header
fields (\texttt{h=}), a hash of the canonicalized body (\texttt{bh=}), the
canonicalization (\texttt{c=}) and an optional body length limit
(\texttt{l=}). The signature covers the selected header fields, each taken
from the bottom of the header block upward, followed by the
\texttt{DKIM-Signature} field itself with its \texttt{b=} value empty. Under
\emph{simple} header canonicalization the fields are signed as they are;
under \emph{relaxed} canonicalization names are lowercased, folded lines are
unfolded, and runs of whitespace become one space. Naming a field in
\texttt{h=} more often than it occurs (``oversigning'') makes later
additions of that field break the signature. RSA keys have at least 1024
bits~\cite{rfc8301}. A DKIM signature attests to the signing domain, not to
the author: DMARC~\cite{rfc7489} calls a signature \emph{aligned} when
\texttt{d=} matches the domain of the \texttt{From} address, whose syntax,
with display names, quoting, comments and folding, is
RFC~5322's~\cite{rfc5322}.

\subsection{Smart-account recovery on Perch}
Soroban is the smart-contract platform of the Stellar
network~\cite{soroban-docs}. A contract account authorizes through its own
\texttt{\_\_check\_auth}, which the host calls when a contract it is invoked
by requires its authorization; the call receives the signature payload, the
account's signature value from the transaction's authorization entry, and
the contexts being authorized. Perch~\cite{perch-repo} builds smart accounts
on OpenZeppelin's Stellar contracts~\cite{oz-stellar-contracts} and specifies
their recovery~\cite{perch-recovery-spec}.

An account's policy document may carry a \emph{recovery member}: a profile,
a mode, a controller, the replaceable signers, and timing in ledgers. Under
the \texttt{Loss} profile the owner key may be lost but not stolen, so the
owner can veto any recovery; under \texttt{Protected} it may be stolen, so
changes to recovery also need the recovery condition. A
\texttt{guardian-only} mode lists guardians and a quorum, \texttt{zk-only}
enrolls one zero-knowledge credential, and \texttt{combined} needs both.
Perch binds the member, guardian addresses included, into a
\texttt{config\_hash}.

Anyone may open a recovery attempt, which grants nothing by itself. Each
attempt binds one \emph{recovery statement}, a fixed-width encoding of the
action (lost-key, compromise, cancel, reconfigure or upgrade), the network,
the account, the controller, the configuration epoch and hash, timing, and
an action-specific subject; for a lost-key recovery the subject names the
attempt id and commits to the replacement credentials and the target
document. Guardians authorize the statement's digest
$\digest = \Sha(\text{statement})$ through the controller's
\texttt{submit\_guardian}, which calls
\texttt{guardian.require\_auth\_for\_args((d,))} and counts each guardian at
most once per attempt and evidence domain (initiation or cancellation);
reconfigurations and upgrades collect approvals of a change digest the same
way through \texttt{approve\_change}. When the quorum is reached the attempt is authorized;
it completes only after the delay, during which a \texttt{Loss} owner may
cancel it. Every configuration change and completion bumps the epoch, which
invalidates older attempts and evidence.

\subsection{Noir, UltraHonk and Soroban verification}
Barretenberg~\cite{barretenberg-docs} proves Noir~\cite{noir-docs} circuits
with UltraHonk, a Plonk-family~\cite{plonk} proof system whose verifier runs
a sumcheck, as in HyperPlonk~\cite{hyperplonk}, under a universal reference
string, so a new circuit needs a new verification key but no setup of its
own. Perch verifies UltraHonk proofs in Barretenberg's zero-knowledge flavor
with a Keccak transcript in a Soroban contract; proofs are 16{,}224 bytes
whatever the circuit's size\claim{circuit-size}, and the email circuit uses
the same verifier code with its own key. Poseidon2~\cite{poseidon2} is the
hash shared by Perch's circuit and Soroban's host.

\section{Threat model and goals}\label{sec:threat}

\paragraph{Parties.} The \emph{user} holds a passkey for the account and,
once it is lost, a new device. The \emph{mail provider} signs the user's
outgoing mail and publishes its keys in DNS. The \emph{DKIM key registry} is
a contract governed by a $k$-of-$n$ multisig and fed by independent DNS
attesters. The wallet's \emph{operator} serves the wallet, relays
transactions, controls the recovery address's domain, and runs the alert
service: an untrusted host and a \emph{sealed vault}, an attested enclave
that alone reads alert addresses and inbound mail. Alerts leave through a
third-party email-sending service. The account's other guardians are
whatever its recovery policy lists.

\paragraph{Adversaries.} We consider eight, alone or together. The
\emph{inbox thief} controls the victim's inbox, Sent folder included, and
wants the account. The \emph{social engineer} can get the user to send an
email with a chosen Subject to a chosen address and to read out or type the
recovery code, for example posing as support. The \emph{registry operator}
is registry members below the threshold, or a threshold of them adding a key
the provider never published. The \emph{relayer} or \emph{front-runner}
drops, delays, reorders or copies the user's transactions, or submits a copy
first, for another statement, attempt or account. The \emph{mail provider}
sees the user's mail and holds the signing key; it wants to learn which
account its user recovers. The \emph{network observer} sees the chain and
every pending transaction. The \emph{malicious frontend} is a misbehaving
wallet build: at recovery an attacker like any other client, at enrollment
trusted (A7). The \emph{operator's staff} read and change the alert host's
process, storage and logs, and the relayer, and want users' addresses; they
do not alter the recovery domain's mail routing (its MX, MTA-STS policy or
mail exchanger's TLS keys) or the wallet build users run (A10, A11), and
against staff who do, nothing built detects it (\S\ref{sec:limits}). Any adversary may open recovery
attempts against any account, send arbitrary email and run its own mail
domain. None breaks the cryptographic assumptions, reads the enclave's
memory, or holds two of the vault key's three shares.

\paragraph{Assumptions.} Table~\ref{tab:assumptions} lists what the
properties rest on; each is a ledger entry with status \claimtag{assumed}.
A1 is split: unforgeability (A1a) is cryptographic; sender authentication
(A1b) is a property of each provider's submission path, which
Appendix~\ref{app:providers} reviews from public sources, with every gap
open; a controlled submission test from real accounts is a check before
mainnet. If A1b fails at a provider, for example by checking only the first
of several \texttt{From} fields and signing the last, the factor there reduces to knowing
the victim's address and $\code$, which still open the commitment, under
the same delay, quorum and alerts\claim{email-factor-soundness}. A9
separates two clocks: the windows this design adds are measured on ledger
close time and keep their wall-clock length whatever the network's ledger
rate; Perch's delay and the statement's expiry are ledger counts.

\begin{table}[t]
  \caption{Assumptions. Each is \claimtag{assumed} in the claim ledger.}
  \label{tab:assumptions}
  \small
  \begin{tabularx}{\columnwidth}{@{}l>{\raggedright\arraybackslash}X@{}}
    \toprule
    Id & Assumption \\
    \midrule
    A1a & RSASSA-PKCS1-v1\_5 with SHA-256 is unforgeable for 2048-bit moduli,
          and providers keep their keys secret\claim{dkim-unforgeability}. \\
    A1b & A provider signs a \texttt{From} naming one of its mailboxes only
          for that mailbox's holder: it refuses two \texttt{From} fields or
          authenticates the one it signs, sending as another in-domain
          address needs proof of control, and a reassigned address is
          useless without $\code$\claim{provider-authenticates-sender}. \\
    A2, A3 & SHA-256 and Poseidon2 are collision resistant, and no known
          input maps to zero under Poseidon2; for hiding, Poseidon2 is a
          random function\claims{hash-collision-resistance,commitment-hiding,blinded-subject-hiding}. \\
    A4 & UltraHonk in Perch's zero-knowledge flavor is zero knowledge and
          weakly simulation extractable for adversaries that see honest
          proofs (stated, not established), and the pinned toolchain compiles
          and proves the circuit as
          written\claims{proof-system-security,toolchain-correctness}. \\
    A5 & Every key the registry lists for a domain was published by that
          domain's provider\claim{registry-correctness}. \\
    A6 & $\code$ and $\blind$ come from a secure generator and stay with the
          user and the browser\claim{secret-generation}. \\
    A7 & The enrolling wallet commits to the real enrollment email's address
          and pins only keys it checked\claim{honest-enrollment}. \\
    A8 & The enclave keeps its memory from its host, and at most one of the
          three key custodians is malicious\claim{enclave-isolation}. \\
    A9 & Close time tracks wall-clock time; ledgers close on average no
          faster than the target close time, from which the wallet sizes
          Perch's 7-day delay floor\claim{ledger-timing}. \\
    A10 & Mail to the recovery address reaches only the enclave: the
          operator does not change the recovery domain's MX, MTA-STS
          policy~\cite{rfc8461} or the mail exchanger's TLS
          keys\claim{recovery-domain-routing}. \\
    A11 & Users run the published wallet build\claim{published-wallet-build}. \\
    A12 & Under email plus delay, the owner learns of an authorized attempt
          within the delay\claim{alert-liveness}. \\
    \bottomrule
  \end{tabularx}
\end{table}

\paragraph{Goals.} The email factor should satisfy four properties, stated
formally in \S\ref{sec:properties}.
\begin{description}
  \item[G1 Soundness.] An accepted approval implies a DKIM signature by a
    registered key for the enrolled domain over an email whose sender is the
    committed address and whose content binds the exact recovery statement:
    its account, new key, nonce and expiry.
  \item[G2 Privacy.] The chain's state, the transactions and the proofs
    reveal nothing about the address beyond the salted commitment.
  \item[G3 Replay and front-running resistance.] An approval is bound to one
    account, one statement and one nonce, and a copy submitted by anyone
    records the same single vote.
  \item[G4 One-factor safety.] No sequence of email approvals alone
    completes a recovery: completion still needs the delay and the policy's
    quorum.
\end{description}
G4 is the precise sense of ``email is never sufficient alone'': it never
skips the delay. Under the lost-key profile with the email guardian as the
only one needed (\emph{email plus delay}: at least seven days, and a second
alert channel registered when the wallet builds the configuration), email
alone authorizes, and what stops a wrongful recovery is the owner's
cancellation within the delay, which rests on A12. The wallet checks the
channel when it builds a configuration; if the channel later lapses or goes
unacknowledged, the chain still accepts the email vote alone, and the wallet
tells the owner and offers to add a device or another helper.

Off chain the goal is scoped. The wallet tells users exactly
\PrivacyClaim, a claim against the operator's servers and staff that holds
under A8, A10 and A11; it is not a claim that the address can never
leak\claim{scoped-privacy-claim}. The user's mail provider and the
alert-sending service necessarily see the address. The provider sees the
blinded Subject, whose content names no attempt; the timing of the email
and of the vote that follows it can (\S\ref{sec:limits}).

\paragraph{Beyond the protocol.} G1 binds the email to the statement, not to
the sender's informed consent. The social engineer needs no fake page: the
real recovery page on its own device opens a lost-key attempt for its key
(opening is permissionless) and shows a Subject the user cannot check, and
since alert codes derive from $\code$, the attacker holding it sees the code
the alerts carry. When the victim acts, typing the dictated Subject into
their own mail client, no text of the operator's reaches them: the prepared
body's warning and the warning where the code is typed render on the
attacker's screen and aid honest users only. The alerts reach the victim
after the vote, and each required alert says that whoever told the reader to
expect it, or to ignore it, is trying to take the
account\claim{phishing-defences}. What remains is the delay with the owner's
cancellation, the push to the device set up for alerts, the wallet's own read
of open attempts, and other guardians; a \texttt{Re:} reply to a phishing
message is unprovable. Whether the recovery address's mail exchanger should
reply to each approval's sender with a warning, with \texttt{To:} bound
where signed so the reply cannot be skipped, is open
(Appendix~\ref{app:open}). Two risks remain under email plus
delay: a phished user who has lost every device cannot cancel, and with
alerts going to a stolen inbox and no push device reachable, only opening the
wallet shows the attempt (\S\ref{sec:eval-attacks}). Whether users act on
what reaches them is not measured\claim{user-behaviour}.

\paragraph{Which losses the factor covers.} Platforms that sync passkeys
sync the wallet's by default, so losing every device loses such a passkey
only with the platform account or its keychain's unlock
secret\claim{passkey-sync}. For a Gmail or iCloud user the platform account
can also hold the factor's inbox (Table~\ref{tab:losses}); enrollment notes
such an inbox and recommends another inbox or
helper\claim{enrollment-safeguards}.

\begin{table}[t]
  \caption{When the email factor adds recovery coverage, by passkey and
    loss, assuming the inbox's provider does not require a lost device to
    sign in.}
  \label{tab:losses}
  \small
  \begin{tabularx}{\columnwidth}{@{}>{\raggedright\arraybackslash}p{0.2\columnwidth}>{\raggedright\arraybackslash}X>{\raggedright\arraybackslash}X>{\raggedright\arraybackslash}X@{}}
    \toprule
    Lost & Device-bound passkey & Synced, inbox at the sync provider & Synced, inbox elsewhere \\
    \midrule
    Every device & covers & not needed: sync restores the passkey & not needed \\
    Platform account & --- & adds nothing: the inbox is lost too & covers \\
    Keychain unlock secret only & --- & covers & covers \\
    \bottomrule
  \end{tabularx}
\end{table}

\paragraph{Non-goals.} The design does not hide that an account has an email
guardian, that a recovery happened, or, from the mail provider, that its user
sent a recovery approval and when. An adversary with the inbox and the
account code who meets the quorum with other factors, or outlasts the delay
unseen, recovers the account: that is the policy's semantics. A compromised
user device, custom mail domains and a provider that signs mail its user did
not send are out of scope.

\section{Design}\label{sec:design}

\begin{figure*}[t]
  \centering
  \resizebox{\textwidth}{!}{\input{figures/sequence}}
  \Description{A message sequence between six parties: the wallet in the user's browser, the user's mail provider, the sealed vault, the Perch recovery controller, the per-account email guardian, and the DKIM key registry with the proof verifier. In enrollment, the wallet draws an account code, the user sends an enrollment email that the provider signs and the vault seals back to the browser, the wallet checks it and computes the commitment and the pin, deploys the guardian, and applies a document listing it. In recovery, an attempt is opened, the controller returns the statement and its digest, the wallet draws rho and computes the blinded value T, the user sends an email whose Subject ends with T, the provider signs it and the vault seals it back, the browser proves the circuit's obligations, the wallet submits the approval, the controller calls the guardian's check, which asks the registry and the verifier, and records one vote. Other guardians reach the quorum, the attempt is authorized, the alert service sends an email and a Web Push with a code, and after the delay, unless the owner cancels, the recovery completes.}
  \caption{Enrollment and recovery. Dashed arrows are answers. The proof is
    made in the browser; the vault, an attested enclave, is the only party
    of the operator's that sees an email or an address.}
  \label{fig:sequence}
\end{figure*}

Figure~\ref{fig:sequence} shows the two flows. Each account that opts in gets
its own immutable email-guardian contract, listed as one guardian in its
Perch recovery member and counted toward the quorum like any other. The
guardian fixes at construction the commitment $\Cm$ to the enrolled address,
the account, its Perch controller, the DKIM key registry it trusts, an
optional key pin, the actions it may approve, and the proof verifier.

\subsection{Statement binding}\label{sec:design-statement}
The email factor approves exactly Perch's recovery statement and nothing
else; there is no email-specific statement. The guardian authorizes
$\digest = \Sha(S)$ for the statement encoding $S$, the circuit binds
$\digest$, and the guardian also receives $S$ itself. The four items a
recovery must bind are fields of $S$:
\begin{description}
  \item[Account] the \texttt{account} contract id, which the guardian
    requires to be its own.
  \item[New key] for lost-key and compromise recovery, the hash of the
    replacement credentials and of the derived target document; for a
    reconfiguration the new configuration hash, for an upgrade the Wasm hash.
  \item[Nonce] the attempt id, which Perch never reuses for an account
    (request id for upgrades), and the configuration epoch, which every
    configuration change and completion bumps.
  \item[Expiry] \texttt{valid\_until\_ledger}, a ledger sequence number;
    Perch's statement carries no wall-clock time, so the DKIM \texttt{t=}
    tag is not used.
\end{description}
The action is inside $\digest$, so one Subject format serves every action,
and the guardian checks the action against its allowed set. The circuit takes
the digest as two 128-bit limbs $\digest_{hi}, \digest_{lo}$, exactly as
Perch's ZK projection splits it, so both are always canonical field elements.

\subsection{The recovery email and its blinding}\label{sec:design-email}
The user sends the recovery email directly from the enrolled inbox to
\textsf{RECOVERY\_INBOX}, one fixed address per deployment; nothing about
the user is needed to start recovery. The wallet opens the user's mail client
with the recipient and Subject filled in. The Subject is exactly
\begin{gather*}
  \texttt{APPROVE NIDO ACCOUNT RECOVERY } \mathrm{hex}(T),\\
  T = \blindf(\digest, \blind) = \Poseidon(\mathsf{DOM\_BLIND}, \digest_{hi}, \digest_{lo}, \blind),
\end{gather*}
where $\blind$ is 31 fresh random bytes that never leave the browser and
$\mathrm{hex}$ gives 64 lowercase hexadecimal characters. Every statement and
its digest are public, so a Subject carrying $\digest$ in the clear would let
the provider match the email to the account and attempt; $T$ hides
$\digest$ from anyone without $\blind$ when Poseidon2 is modelled as a random
function\claim{blinded-subject-hiding}. Nothing may precede or follow the
value: a \texttt{Re:} prefix, an RFC~2047 encoded word~\cite{rfc2047} or
trailing text makes the email unprovable, while a provider that folds the
long line at a space is harmless, because relaxed canonicalization restores
the exact bytes\claim{subject-blinding}.

The proof binds neither the recipient, nor the body, nor \texttt{t=}: not
every launch provider signs \texttt{To:} or sets \texttt{t=}, an email to any
recipient is equally good evidence of inbox control, and freshness comes from
the statement's deadline and single-use attempt id. The signature still
covers the whole signed header block.

The signed octets reach the browser through the sealed vault, which runs
the recovery address's mail exchanger: the browser first registers $T$ with
a fresh device key, sealed to the vault with HPKE~\cite{rfc9180}; the first
registration of a value holds it; and the vault seals each matching message
to that key and keeps nothing\claim{sealed-inbound}. Alternatively the user
imports the sent message as an \texttt{.eml} file.

\subsection{Enrollment and the commitment}\label{sec:design-enroll}
While the owner still holds the passkey, the wallet draws $\code$, 16 bytes
from a secure generator, never reused and never sent to Nido. The user sends
an enrollment email whose Subject is a distinct phrase followed by
$\mathrm{hex}(E)$, where $E$ hashes the account's contract id under its own
tag with a fresh $\blind_e$, and the wallet reads the exact
\texttt{From} address $\addr$ from the DKIM-signed copy, so the user never
types it and the provider writes the same bytes on every later email. The
commitment is
\begin{align*}
  \Cm = \commitf(\addr, \code) = \Poseidon(&\mathsf{DOM\_COMMIT}, |\addr|,\\
    &\mathrm{pack}(\addr), \code),
\end{align*}
where $\mathrm{pack}$ cuts the address, right-padded to 124 bytes, into four
31-byte field elements and each $\mathsf{DOM}$ is a distinct domain tag; the domain of the address is inside $\Cm$, so the
enrolled domain is the committed address's. The wallet checks the enrollment
email under the recovery circuit's own rules before deploying a guardian, so
an inbox whose recovery email could never be proven is refused at
enrollment\claim{enrollment-email-check}; an enrollment Subject can never
satisfy the recovery circuit\claim{enrollment-subject-distinct}. Nothing on
chain verifies the enrollment email: a wrong $\Cm$ disables only its own
owner's email factor.

The wallet then deploys the guardian with $\Cm$ and a key pin
(\S\ref{sec:design-registry}), reads back its configuration and code hash,
and applies a document that lists it. Listing is an ordinary Perch
reconfiguration, so under \texttt{Protected} it needs the current recovery
condition too; since the guardian is immutable, changing the address, the
code, the pin or the allowed actions means deploying a new guardian through
the same gate.

Recovery needs $\code$, so the wallet offers two backups: a block in the
enrollment email's body, kept in the Sent folder, encrypting $\code$ under an
Argon2id~\cite{rfc9106} key from a wallet-generated five-word passphrase; and
a recovery kit\claim{account-code-backup}.

\subsection{The circuit relation}\label{sec:design-circuit}
The recovery circuit (Fig.~\ref{fig:relation}) has four public inputs, in
this order: $\Cm$, $\digest_{hi}$, $\digest_{lo}$, and $\mathit{root}$, the
root of the key tree the signing key is proven against. The guardian builds
all four itself. The private inputs are the signed header block $H$ (at most
1536 bytes), the RSA signature $s$ and modulus $n$, $\blind$, $\code$, and the
key's leaf index and 16 Merkle siblings. Positions inside $H$ are hints that
the circuit computes in unconstrained code and then checks. There is no
action input, timestamp or nullifier: the action and the expiry are in the
digest, and Perch's controller, not a nullifier, stops double counting.

\begin{figure}[t]
  \centering
  \resizebox{\columnwidth}{!}{\input{figures/circuit-relation}}
  \Description{Three columns. Private inputs on the left: the signed header block H of at most 1536 bytes, the RSA signature and modulus, the account code, rho, and the leaf index with 16 siblings. Obligations C1 to C6 in the middle. Public inputs on the right: the commitment, the two digest limbs, and the key-tree root. Arrows show that the header feeds C1 to C5, the signature and modulus feed C1 and C6, the account code feeds C3, rho feeds C5, the Merkle path feeds C6, and C3, C5 and C6 open the commitment, the digest limbs and the root.}
  \caption{The recovery circuit's relation. Domain tags and lengths inside
    the Poseidon2 hashes are omitted.}
  \label{fig:relation}
\end{figure}

The circuit enforces six obligations:
\begin{description}
  \item[C1] $s<n$, $n$ has exactly 2048 bits, and $s^{65537} \bmod n$ is the
    EMSA-PKCS1-v1\_5 encoding of $\Sha(H)$.
  \item[C2] $H$ is a well-formed signed block: header fields in relaxed
    canonical form, each ended by CRLF, then the signature's own field with
    an empty \texttt{b=} and no trailing CRLF, carrying
    \texttt{v=1}, \texttt{a=rsa-sha256}, relaxed header canonicalization,
    an \texttt{h=} naming \texttt{from} and \texttt{subject}, and a
    \texttt{d=}; no byte past $|H|$ is read.
  \item[C3] $H$ holds exactly one \texttt{from} field; its value parses to an
    address $\addr$ that, with $\code$, opens $\Cm$.
  \item[C4] The lowercased domain of $\addr$ equals the lowercased
    \texttt{d=}, which is the domain of the registry leaf.
  \item[C5] $H$ holds exactly one \texttt{subject} field, whose value is the
    recovery phrase followed by $\mathrm{hex}(T)$ for the public digest limbs
    and $\blind$.
  \item[C6] The leaf $\leaff(\mathit{dom}, \kh(n))$ has a depth-16 Merkle
    path to $\mathit{root}$, where $\kh(n) = \Poseidon(\mathsf{DOM\_KEY},
    n_0,\ldots,n_{17})$ over the 18 120-bit limbs of $n$, and the leaf hashes
    $\mathsf{DOM\_LEAF}$, the domain's length, the domain packed into two
    field elements, and $\kh(n)$.
\end{description}
The circuit takes the canonicalized bytes the signature covers and never
canonicalizes; a field outside $H$ has no effect. The \texttt{From} value
must be a bare address or a printable display name followed by exactly one
angle-bracketed address that ends the field; quoted local parts, comments,
groups and address lists are refused, the address has at most 124 bytes,
and it is used exactly as signed, with no case folding or tag
removal\claim{circuit-from-grammar}. The circuit refuses simple header
canonicalization, a missing or duplicated \texttt{From} or \texttt{Subject},
a misaligned \texttt{d=}, and a signature not below its
modulus\claim{circuit-dkim-rules}. Changing any public input, or any private
value that links the email to one, makes the proof
fail\claim{circuit-binds-public-inputs}, and an honest recovery email from
every launch domain is accepted, including a 124-byte address with a 500-byte
display name\claim{circuit-accepts-launch-providers}.

\subsection{The email guardian}\label{sec:design-guardian}
The guardian's constructor refuses an account or controller that is not a
contract, the account as its own controller, a zero or non-canonical
commitment, a pin that has expired or lasts more than 365 days, an empty or
unknown action set, and a verifier whose reported circuit id differs from the
configured one\claim{guardian-construction}. Perch's controller calls
\texttt{guardian.require\_auth\_for\_args((d,))}, and the host then runs the
guardian's \texttt{\_\_check\_auth} with the approval $(S, \mathit{root},
\pi)$ from the guardian's authorization entry. The guardian refuses at the
first failing check of Table~\ref{tab:guardian}, and builds the public
inputs only from its configuration, the authorized digest and the accepted
root\claim{guardian-checks}. Approvals by the email guardian count toward a
\texttt{Protected} reconfiguration through \texttt{approve\_change} like
any guardian's\claim{guardian-on-perch}.

\begin{table}[t]
  \caption{The guardian's checks, in order.}
  \label{tab:guardian}
  \small
  \begin{tabularx}{\columnwidth}{@{}r>{\raggedright\arraybackslash}X l@{}}
    \toprule
    \# & Check & Error \\
    \midrule
    1 & One contract call, on the controller, with one 32-byte argument $\digest$ & \texttt{UnexpectedContext} \\
    2 & $\Sha(S) = \digest$ & \texttt{DigestMismatch} \\
    3 & $S$ is a version-2 Perch statement of a known action & \texttt{MalformedStatement} \\
    4--6 & $S$ names this network, account and controller & \texttt{Wrong\ldots} \\
    7 & The action is allowed & \texttt{ActionNotAllowed} \\
    8 & The call is that action's entry point & \texttt{WrongEntryPoint} \\
    9 & Ledger $\le$ \texttt{valid\_until\_ledger} & \texttt{StatementExpired} \\
    10 & $\mathit{root}$ is canonical & \texttt{NonCanonicalRoot} \\
    11 & $\mathit{root}$ is the unexpired pin, or the registry accepts it & \texttt{UnknownRoot} \\
    12 & $\pi$ verifies for $(\Cm, \digest_{hi}, \digest_{lo}, \mathit{root})$ & \texttt{ProofRejected} \\
    \bottomrule
  \end{tabularx}
\end{table}

\paragraph{One guardian per account.} Perch binds guardian addresses into
\texttt{config\_hash}, so everything fixed at the guardian's address is
covered by Perch's reconfiguration gate. A guardian shared by many accounts
would need a table from account to commitment, and whoever could change that
table could change an account's email factor without passing the gate. And
\texttt{require\_auth\_for\_args} hands the guardian only $\digest$: it
learns which account $\digest$ concerns only from $S$ (check~2), and approves
only for its own (check~5).

\paragraph{What the guardian does not check.} It ignores the host's
signature payload (the entry's nonce, expiration and invocation hash): the
proof binds the statement, not an authorization entry, so anyone can re-wrap
an approval, and the re-wrapped copy can still authorize only the same
digest, which the controller counts once\claim{guardian-payload-irrelevant}.
Quorum, delay, cancellation and epoch are the controller's; the guardian
could not read them anyway, since Soroban refuses re-entry into the
controller that is calling it. The guardian calls only the registry and the
verifier, which embeds Perch's UltraHonk verifier with the circuit's
verification key\claim{verifier-contract}.

\subsection{The DKIM key registry and the pin}\label{sec:design-registry}
Keys are identified by $\kh(n)$, never by selector, because a provider may
reuse a selector for a new key. The registry holds a depth-16 Merkle tree of
leaves $\leaff(\mathit{domain}, \kh)$ hashed exactly as C6 hashes them, with the
host's Poseidon2, and answers one question for the guardian: is this root
accepted now\claim{registry-hash-parity}? Its rules are:
\begin{itemize}
  \item \emph{Activation.} An addition, proposed and approved by the
    multisig, executes 72 hours after the later of its approval and the
    ledger since which a threshold of attesters' latest observations see the
    key live in DNS, or 7 days after its approval, whichever comes
    first\claim{registry-timelocks}.
  \item \emph{Grace.} A root superseded at ledger $L$ stays acceptable
    through $L + \mathit{grace}$, so a proof against the previous root still
    verifies and a key the provider has just retired keeps working briefly.
  \item \emph{Revocation.} An emergency threshold of members revokes a key
    for compromise; from that ledger every root containing it is refused,
    grace included, a pending addition of it stops, and it can return only
    through a new full-threshold addition on the 7-day
    path\claim{registry-revocation}.
  \item \emph{Scope.} Leaves exist only for allowlisted domains and RSA keys
    of exactly 2048 bits; the registry computes key hashes itself; the empty
    tree's root is never accepted\claim{registry-governance}.
\end{itemize}
An attester resolves \texttt{selector.\_domainkey.domain}, signs an exact
encoding of what it saw with Ed25519~\cite{rfc8032}, bound to the registry's
contract id, and anyone may submit the signature; each attester's latest
observation counts\claim{registry-attestations}. The registry has no admin,
setter or upgrade path\claim{registry-immutable}: changing a parameter
means deploying a new registry, which a guardian follows only by being
replaced through Perch's gate. Its proposed governance is five members with
threshold three, an emergency threshold of two, three attesters with
threshold two, and a 24-hour grace period\claim{registry-defaults}.

\paragraph{The per-account pin.} A guardian may carry one pin, a key tree
fixed at construction, with an expiry of at most 365 days. Its leaf 0 is
$\Poseidon(\mathsf{DOM\_PIN\_SALT}, \code)$ and its other leaves are
registry-format leaves of the pinned keys, so the pin's root reveals nothing
about which provider's keys it holds without $\code$. The guardian accepts a
key root if it is the unexpired pin's or the registry accepts it, so the pin
carries a new account through the window while its provider's key waits in
the registry's timelock, and keeps the enrolled key usable if the registry
drops it\claim{key-pin}. The pin is the one path the registry cannot close: a
revoked key stays usable through a pin that holds it until the pin expires,
and forging an approval with it still needs the victim's address and
$\code$. The wallet deploys guardians with a 30-day
pin\claim{wallet-pin-lifetime}.

\paragraph{Finding the key.} Asking for one provider's key reveals the
provider, so the wallet reads a static list of every registry key, the same
for every user, and falls back to DNS over HTTPS only after telling the
user; either way a key counts only through the pin or an accepted
root\claim{key-lookup}.

\subsection{Alerts and the sealed vault}\label{sec:design-alerts}
Alerts are driven by the controller's own events; nothing is written on
chain for them. The alert service's host watches the controller by contract
id and hands each event to the vault, which finds the account's record under
$\mathrm{HMAC}(k_\mathit{alert}, \mathit{account})$ with a key that never
leaves it. Authorization, completion and configuration changes alert every
channel; opening an attempt, which anyone can do, at most one email a day.
An alert has generic text and no link, and carries a six-digit code keyed by
a key derived with HKDF~\cite{rfc5869} from $\code$, which the wallet shows
for the attempt it knows\claim{alert-code}.

When email is a recovery factor, a second channel is required: Web Push to a
passkey device~\cite{rfc8030,rfc8291,rfc8292}. A contact is registered,
replaced or removed only with an assertion by one of the account's
administrative passkeys~\cite{webauthn3}; a replacement or removal takes
effect after 14 days, longer than the 7-day minimum delay, while the old
contact keeps receiving alerts and is told at once; and removal is refused
while a recovery is authorized\claim{alert-contacts}. Alert email goes
through an EU-owned sending service, with failover to a second
one\claim{alert-formats}.

The browser encrypts each address to the vault with HPKE after verifying the
vault's attestation document~\cite{aws-nitro-attestation}. The vault's master
key is split 2-of-3 with Shamir's scheme~\cite{shamir} across independent
custodians, who release a share only to an enclave image that an on-chain
image registry has activated after a 14-day public timelock, and stop once
the image is revoked\claim{vault-custody}. With every third party mocked and
the vault in its own process, no message, host record, log, RPC call,
custodian decision or push body contains the address: Nido's servers and
staff can't read your address\claim{servers-cannot-read-address}, under A8.
The sending service and the user's mail provider necessarily see it.

\subsection{Wallet policy}\label{sec:design-policy}
Perch's controller accepts any configuration its owner applies, so the
remaining rules bind Nido's wallet, which builds every recovery member
through one builder that tells the email guardian apart by a typed marker,
never by address:
\begin{enumerate}
  \item at most one email guardian per recovery
    member\claim{wallet-one-email-guardian};
  \item an email guardian only with a registered second alert
    channel\claim{wallet-second-alert-channel};
  \item if the email vote alone can authorize (guardian-only, quorum one),
    the \texttt{Loss} profile and a delay of at least seven days, 120{,}960
    ledgers\claim{wallet-email-plus-delay};
  \item every action allowed whenever the quorum needs the email vote, so
    that a \texttt{Protected} account can still reconfigure and
    upgrade\claim{wallet-email-allowed-actions}.
\end{enumerate}
OIDC sign-in with Google, Apple or Microsoft may later become a second way
to produce the email factor; it would never count as a factor independent of
the same provider's DKIM, since both rest on one provider.

\section{Security analysis}\label{sec:properties}

Each property of \S\ref{sec:threat} is a theorem of a Lean~4
model~\cite{lean4}. The model contains the guardian's \texttt{\_\_check\_auth}
as the function $\mathsf{check}$, its twelve checks in order with their error
codes; the circuit's relation as $\mathsf{Relation}(x, w)$, one field per
obligation C1 to C6 for public inputs $x = (\Cm, \digest_{hi}, \digest_{lo},
\mathit{root})$ and witness $w$; and Perch's controller for one account as a
state machine whose transitions cite Perch's transitions T1 to T9, with an
event log that records the configuration each event happened under.
$\mathsf{check}$ succeeds exactly when the facts of its twelve checks
hold\claim{check-iff-accepted}. SHA-256 and Poseidon2 are abstract
functions, and their breaks (a collision, or a Poseidon2 input that maps to
the empty leaf $0$) are named propositions. The header-level functions of
the circuit (which bytes are the \texttt{From} address, the Subject value,
the \texttt{d=} value; lowercasing; hex; the PKCS\#1 encoding) are abstract
too: computing them as \S\ref{sec:design-circuit} says is the circuit's
obligation, tested in \S\ref{sec:hardening} but not proven. No assumption is
an axiom: the build pins every theorem's axiom dependencies to Lean's
built-in three\claim{lean-axioms}.

\subsection{Soundness}

\begin{theorem}[Soundness of the email factor]\label{thm:soundness}
Let guardian $G$ have configuration $c$, and suppose
$\mathsf{check}(c, \mathit{env}, (S, R, \pi), \mathit{ctx})$ succeeds at
ledger $L$. If the configured verifier has a knowledge extractor (A4), there
are an entry point $f$, a digest $\digest$ and a witness $w$ such that:
\begin{enumerate}
  \item $\mathit{ctx} = [\,c.\mathit{controller}.f((\digest,))\,]$,
    $\Sha(S) = \digest$, and $S$ names $c$'s account, controller and
    network, an allowed action whose entry point is $f$, and a
    \texttt{valid\_until\_ledger} $\ge L$;
  \item $R$ is the pin's root with $L$ before its expiry, or a root the
    registry accepts at $L$;
  \item $w.H$ is well formed and carries a valid RSA-2048 PKCS\#1 v1.5
    signature $w.s$ under $w.n$ (C1, C2);
  \item $\leaff(w.\mathit{dom}, \kh(w.n))$ has a Merkle path to $R$ (C6);
  \item the \texttt{From} address $a$ of $w.H$ has at most 124 bytes and
    $\commitf(a, w.\code) = c.\Cm$ (C3);
  \item $\mathrm{lower}(\mathrm{domain}(a)) = \mathrm{lower}(\texttt{d=}) =
    w.\mathit{dom}$ (C4);
  \item the Subject of $w.H$ is the recovery phrase followed by
    $\mathrm{hex}(\blindf(\digest, w.\blind))$ (C5).
\end{enumerate}
If, moreover, every root the registry accepts is the root of a depth-16 tree
whose non-empty leaves are trusted (A5), and $c.\Cm$ commits to the enrolled
address $a^\ast$ while the pin holds only trusted leaves (A7), then
$a = a^\ast$ and the key's leaf is trusted, or Poseidon2 has a collision, or
Poseidon2 maps some input to $0$.
\end{theorem}
The first part is \textsf{EmailRecovery.email\_of\_check} and the whole is
\textsf{EmailRecovery.soundness}\claim{email-factor-soundness}. Through $S$,
the email binds the four items of \S\ref{sec:design-statement}. The step to
G1 is A1: a valid RSA-2048 signature under a trusted key of
$\mathrm{domain}(a^\ast)$ over a header whose \texttt{From} is $a^\ast$ was
produced by that domain's provider, which signed it for the holder of
$a^\ast$. That step is computational and outside the model. The Poseidon2
breaks matter because the key trees' empty leaf is $0$: a key leaf that
hashed to $0$ would stand in for an empty slot of any tree, which collision
resistance alone does not exclude, so A3 includes it. The proof uses two
lemmas about the spec's formulas: the commitment binds the address and the
code, and a Merkle path binds its leaf, each unless Poseidon2 collides.

\subsection{Privacy}

\begin{theorem}[Privacy as noninterference]\label{thm:privacy}
Let a guardian be enrolled from a secret $(a, \code, K)$, an address, an
account code and pinned keys. There is a function $f$ such that, for every
initial controller state $s_0$ and every sequence of transactions, the
chain's view of the run (the guardian's configuration, every verdict it
gives, and the controller's state and event log) equals
\[
  f\bigl(\commitf(a, \code),\ \mathrm{pinRoot}(\code, K),\ s_0,\ \text{transactions}\bigr).
\]
\end{theorem}
This is \textsf{EmailRecovery.privacy\_noninterference}, with the two-run form
\textsf{privacy\_two\_runs} and \textsf{publicInputs\_enroll}, which shows
the proof's public inputs are $\Cm$, the digest limbs and the
root\claim{address-privacy}. The theorem reduces privacy to the secrecy of
the preimages of $\Cm$ and the pin's root. G2, the computational statement
that a guardian's state and accepted approvals are indistinguishable between
two addresses at allowlisted providers when $\code$ is uniform and unknown,
follows only with A3's random-function model, A4's zero knowledge and A6,
and is not proven\claim{commitment-hiding}\claim{secret-generation}. What
remains visible is that the account has an email guardian, its allowed
actions, when each approval happened, and whether it used the pin or a
registry root. Off chain the provider sees $T$, which hides $\digest$ under
the same assumptions\claim{blinded-subject-hiding}.

\subsection{Replay and front-running resistance}

\begin{theorem}[Replay resistance]\label{thm:replay}
(a) If $\mathsf{check}$ accepts $(S, R, \pi)$ in contexts $\mathit{ctx}$, then
$\mathit{ctx} = [\,c.\mathit{controller}.f((\Sha(S),))\,]$; it accepts it in
exactly one context and for exactly one account, controller and network.
(b) Two statements with equal digests have the same account, network,
controller, action, attempt or request id and epoch, or SHA-256 collides.
(c) If the composed controller and guardian record the same approval in two
successful \texttt{submit\_guardian} calls, both are for one attempt, one
evidence domain and one epoch, or SHA-256 collides; a vote already recorded
is refused again. (d) The verdict ignores the authorization entry the
approval is wrapped in, and the controller's transitions take no submitter
input.
\end{theorem}
In Lean, parts (a) to (d) are \textsf{accepted\_names\_digest},
\textsf{accepted\_context\_unique} and
\textsf{accepted\_account\_unique};
\textsf{digest\_fixes}\claim{digest-fixes-statement};
\textsf{Perch.email\_vote\_replay}\claim{replay-resistance} with
\textsf{votes\_counted\_once} and \textsf{counted\_vote\_refused\_again}; and
\textsf{check\_payload\_irrelevant}\claim{guardian-payload-irrelevant} with
\textsf{emailAuth\_payload\_irrelevant}. A copy of an approval submitted by
anyone therefore records the vote the original would have recorded, for the
statement the original named. Part (b) uses two facts about Perch's encoding
as hypotheses: the encoding is injective, and the guardian's parser reads
back every field. They are tested, not proven: Perch's encoder reproduces
independent vectors, and the guardian parses every statement shape it
produces\claim{model-code-agreement}.

\subsection{One-factor safety}

\begin{theorem}[One-factor safety]\label{thm:onefactor}
In every state reachable from an unenrolled account by transactions whose only
recovery evidence is approvals by $G$ (owner-authorized transactions
allowed), every completed recovery of attempt $i$ is preceded in the log by
the authorization of $i$ at a ledger $L_a$ under a configuration $c$ with
mode \texttt{guardian-only}, quorum $1$ and $G$ among its guardians, and
$0 < c.\mathit{delay}$ and $L_a + c.\mathit{delay} \le L_\mathit{complete}$.
\end{theorem}
This is \textsf{EmailRecovery.Perch.one\_factor\_safety}\claim{one-factor-safety}.
It combines a statement about every trace, whatever its evidence (every
completion waits out a positive delay after its
authorization\claim{completion-waits-for-delay}), with one about traces whose
only evidence is $G$'s (every authorization is at quorum one, guardian-only);
so under a quorum of at least two, or in \texttt{combined} or
\texttt{zk-only} mode, email approvals alone authorize
nothing\claim{email-only-never-authorizes}. The same holds for the recovery
condition of a \texttt{Protected} reconfiguration or
upgrade\claim{email-alone-change-condition}, and during the delay a
\texttt{Loss} owner's cancellation of a live attempt
succeeds\claim{owner-can-cancel}. Under the wallet policy of
\S\ref{sec:design-policy}, stated as a predicate on recovery members, the
email guardian alone reaches the quorum only under \texttt{Loss}, with a
delay of at least 120{,}960 ledgers and a second alert channel
(\textsf{one\_factor\_under\_wallet\_policy})\claim{wallet-policy-one-factor}.
The chain accepts any configuration its owner applies, so this last part
binds Nido's wallet, not other wallets.

\subsection{Obligations}

Table~\ref{tab:obligations} lists, per property, what each component must
do and how it is established. The circuit's obligations and the contracts'
checks are tested; the composition is proven over the model; the
cryptographic steps are assumed.

\begin{table*}[t]
  \caption{Obligations of each property and their status.}
  \label{tab:obligations}
  \small
  \begin{tabularx}{\textwidth}{@{}l>{\raggedright\arraybackslash}p{0.12\textwidth}>{\raggedright\arraybackslash}X>{\raggedright\arraybackslash}p{0.17\textwidth}l>{\raggedright\arraybackslash}p{0.25\textwidth}@{}}
    \toprule
    Property & Circuit & Guardian and registry & Perch & Assumptions & Status \\
    \midrule
    1 Soundness & C1--C6 & checks 1--12; construction's circuit-id check; registry activation, grace, revocation and scope, with C6's hashing & statement built from controller state & A1, A2, A3, A4, A5, A7 & composition proven; C1--C6 and checks tested; registry rules tested; hashing cross-checked \\
    2 Privacy & public inputs exactly $(\Cm, \digest_{hi}, \digest_{lo}, \mathit{root})$ & stores $\Cm$ and a salted pin root; emits nothing & --- & A3, A4, A6 & noninterference proven; indistinguishability assumed \\
    3 Replay & C5 binds $\digest$ & checks 1, 2, 5, 6, 8, 12 & one vote per guardian per attempt and domain or per change digest; ids never reused; epochs grow & A2, A3, A4 & proven, with the encoding facts tested \\
    4 One factor & --- & the guardian is one address & quorum, delay, cancellation & --- & proven over all traces; wallet policy proven as a predicate and tested in code \\
    \bottomrule
  \end{tabularx}
\end{table*}

\subsection{Hardening against public issue classes}\label{sec:hardening}
Public audits of zkEmail's Circom and Noir
libraries~\cite{audit-zksecurity-2024,audit-yacademy,audit-zellic-2024,audit-matterlabs-2024,audit-veridise-2024,audit-consensys-2024}
and public issues in its repositories report recurring classes of
soundness bugs, mostly in header parsing rather than in DKIM itself; such
under-constrained circuits are a common class of SNARK
vulnerability~\cite{chaliasos-sok-snark}. The rules of
\S\ref{sec:design-circuit} refuse each publicly documented class;
Table~\ref{tab:hardening} in Appendix~\ref{app:hardening} lists the ten
classes, their public sources and the refusals. Each class has a negative test that names its source;
the test signs its header with the test key, so the RSA check passes and the
refusal comes from the check under test, and pins the exact failing
assertion\claim{circuit-issue-classes}; an attack on prover hints starts from
the honest hints and changes them as an attacker would. Three classes from the same reports do not
arise: partial-SHA state aliasing (no partial hashing is used), an unbound
Barrett parameter (it enters no constraint and is not hashed), and an
all-zero SHA output (a Circom template not used here).
Beyond the per-class tests, every adversarial email of the provider mocks,
given the witness most favourable to its claim, is refused at a named
assertion whenever its claim must not be provable\claim{circuit-refuses-attacks}
(Appendix~\ref{app:circuit}). Every unconstrained value is a hint that
constrained code re-checks, so a wrong hint can make an honest proof fail
but not make a false statement provable; the pinned compiler's public
advisories were reviewed against the constructs the circuit uses, and its
correctness beyond that review is assumed\claim{toolchain-correctness}.

\subsection{From the model to the code}\label{sec:gap}
The theorems are about the model; the model's agreement with the code is
tested, not proven. The model's $\mathsf{check}$ mirrors the guardian's
Rust \texttt{check} check for check, and each refusal is pinned by a unit
test\claim{guardian-checks}. The controller model mirrors Perch's
\texttt{contract.rs} transition by transition, and the guardian on Perch's
deployed WebAssembly is pinned by integration tests: one approval is one
vote, a resubmission is refused, approvals do not transfer to another attempt
or account, and email approvals alone never complete a
recovery\claim{guardian-on-perch}. The model's abstractions only add
behaviours, so a safety property of the model holds of the contracts:
ledger arithmetic is unbounded, storage never expires, and the target
derivation, the ZK adapter's verdict and the host's acceptance of an
authorization entry are inputs of a transaction. Two identifications remain:
configurations are compared by value where Perch compares
\texttt{config\_hash} (equal unless SHA-256 collides), and the two halves of
a completion are one transition. Not modelled are the account's own state
beyond its applied document, Perch's ZK membership pool, the circuit's
header parsing, and the semantics of Noir and Rust. The wallet-policy
predicate mirrors the wallet's builder rule for rule, and the builder is
tested at each rule's boundary. The registry's tree shape is tested and its
hashing cross-checked against an independent Noir
implementation\claim{registry-hash-parity}.

\section{Implementation}\label{sec:impl}

The implementation is one repository; Perch is a pinned dependency used
unchanged, and the circuit is compiled and proven with exactly Perch's pinned
toolchain, nargo 1.0.0-beta.9 and Barretenberg 0.87.0, so its proofs verify
with Perch's Soroban verifier code\claim{circuit-size}. The optimized
WebAssembly is 33{,}479 bytes for the guardian, 68{,}269 for the registry,
45{,}616 for the verifier and 10{,}152 for the enclave-image
registry\claim{contract-sizes}.

\paragraph{Circuit.} The circuit is a library of components (RSA-2048
PKCS\#1 v1.5 with SHA-256; the structural scan of the signed block; the
\texttt{From}, \texttt{DKIM-Signature} and \texttt{Subject} parsers; packing;
the registry leaf and Merkle root) composed by one relation function. It
vendors only noir-bignum, for the 2048-bit modular multiplication, and
Noir's SHA-256 library, each pinned file by file; Poseidon2 is Perch's
vendored copy; no parser is reused\claim{circuit-vendoring}. The header parsers are written against
the grammars of \S\ref{sec:design-circuit}, and every position in the header
is an unconstrained hint that constrained code re-checks. At a 1024-byte
header bound the same relation has 189{,}503 gates and the same padded size;
the 1536-byte bound costs the prover more but leaves the verification key
and the on-chain cost unchanged\claim{circuit-size}. That bound admits every
launch domain's recovery email and, with a 124-byte address, a display name
of 536 bytes at the tightest provider\claim{header-bound-fits}.

\paragraph{Contracts.} The guardian, the registry, the verifier and the
enclave-image registry are Soroban contracts in Rust. The verifier is a
constructorless, immutable contract that embeds Perch's UltraHonk verifier
from the same Perch commit as the other Perch dependencies, with the email
circuit's verification key compiled in, and reports the key's SHA-256 as its
circuit id; the guardian refuses at construction a verifier whose circuit id
differs from its own configuration\claim{guardian-construction}. A real
proof of each launch domain's recovery email verifies on chain through it,
is refused with any public input changed, and counts as exactly one guardian
vote\claim{onchain-email-vote}. The
registry stores its tree with the host's Poseidon2 sponge and Perch's table
of empty-subtree hashes, so the root it stores is the root the circuit proves
against\claim{registry-hash-parity}.

\paragraph{Wallet.} The browser side is TypeScript: an encoder from a raw
message to the witness, the prover worker (noir\_js 1.0.0-beta.9 and bb.js
0.87.0), which checks its downloads against pinned digests and verifies each
proof locally against the guardian's circuit id\claim{browser-prover}, the
enrollment and recovery pages, and the key lookup.

\paragraph{Alert service.} The host watches the controller over Stellar
RPC with a stored cursor, retries failed alerts, and never repeats one; the
vault runs as a separate process standing in for the enclave. The wallet's
HPKE reproduces RFC~9180's test vector, and the vault's Web Push encryption
RFC~8291's example\claim{alert-formats}.

\paragraph{Mocks.} Every third party is a test-only mock in its published
format: the mail providers sign with test-only keys in each provider's layout,
established from public mailing-list archives\claim{provider-mock-formats};
DNS answers follow RFC~6376's key records and the DNS-over-HTTPS JSON
format\claim{key-lookup}; the sending services follow their published APIs,
Web Push its RFCs, and the enclave's attestation AWS's documented
format~\cite{aws-nitro-attestation}\claim{vault-custody}; the DNS attester
signs test-only Ed25519 attestations.

\section{Evaluation}\label{sec:eval}

The evaluation runs on a local Stellar network and on testnet. It measures
proving time and memory on desktop and phone-class devices; the on-chain cost
of verification in instructions, bytes and fees; and an attack suite covering
forged and replayed emails, a wrong account or key in the statement, stale or
malicious registry keys, a compromised inbox, provider header quirks, and
front-running.

\pending{no measurements have been taken.}

\section{Related work}\label{sec:related}

\paragraph{Proofs about email.} The zkEmail
project~\cite{zkemail-docs,zkemail-verify,zkemail-nr} proves statements
about DKIM-signed email in zero knowledge, with Circom circuits and a Noir
library, and its libraries have been audited
publicly~\cite{audit-zksecurity-2024,audit-yacademy,audit-veridise-2024,audit-consensys-2024}.
This work uses email as one recovery factor inside a delayed, quorum-based
recovery, binds the proof to a recovery statement rather than to an
application command, and blinds that statement from the provider.

\paragraph{Proofs about existing credentials.} zkLogin~\cite{zklogin}
authenticates blockchain transactions with OpenID Connect tokens from
existing platforms and hides the link between the off-chain and on-chain
identity. An OIDC path here would be a second way to produce the email
factor, never an independent one alongside the same provider's DKIM.

\section{Limitations and open questions}\label{sec:limits}

\paragraph{Not built, not measured.} The production enclave is not built:
the binding to the Nitro secure module, TLS terminating inside it, the
recovery address's mail exchanger inside it with the domain's MTA-STS
policy, and a reproducible image build. The custodians, registry members and
attesters are not named. Every third party is a mock: before mainnet each
provider's outbound recovery mail must be checked from a real account,
Fastmail first, whose archived signatures verify under no known key, and a
submission test must settle A1b per provider. Phones are not
measured\claim{phone-proving-budget}, nor registry-root votes on a network,
and no user study was run. There is no external audit, and mainnet waits for
one.

\paragraph{Toward a malicious operator.} \PrivacyClaim{} rests on A8, A10 and
A11, and holds against staff confined to the host (\S\ref{sec:threat}): an
operator that repoints the recovery domain's MX or relaxes its MTA-STS
policy, or serves a modified wallet, reads addresses, and nothing built
detects either\claims{recovery-domain-routing,published-wallet-build}. Three
unbuilt options would make such a change public, and which to build is
open: the domain's MX, MTA-STS policy and TLS key committed in the
enclave-image registry under its 14-day timelock; certificate-transparency
monitoring of the mail exchanger's name; reproducible wallet builds.

\paragraph{Inherent leakage.} The provider learns that its user sent a
recovery approval, and when. The blinded Subject hides which attempt, but
the vote lands seconds later: the browser proves in 2.81\,s on a desktop and
testnet confirmed the vote 1{,}927\,ms after
simulation\claims{browser-proving-cost,testnet-email-recovery}. A provider
watching the chain links the email to its vote, and so to the account,
whenever few email votes land in that window, which can be one; likewise an
enrollment email and the guardian deployment. Unlinkability from the
provider is content-level only. The sending service sees alert addresses and
codes; the alert host learns which accounts register contacts.

\paragraph{Residual risks and burdens.} Under email plus delay, a phished
user with no device left cannot cancel, and an owner whose push device
lapsed or went silent while alerts go to a stolen inbox learns of an attempt
only by opening the wallet; a second helper closes both. The factor is the
inbox and a 128-bit code, so a lost code loses it. It covers the losses of
Table~\ref{tab:losses}; a provider that requires the lost device to sign in,
or a new phone without a mail app, defeats it. A circuit fix does not reach
accounts that cannot reconfigure (\S\ref{sec:design-guardian}), and a
lowered target close time shortens existing email-plus-delay members until
they are rebuilt (A9).

\paragraph{Open questions.} Appendix~\ref{app:open} lists every open
decision with its alternatives, among them a reply to approvals with
\texttt{To:} bound, circuit fixes, the alert code's key, and who pays
restorations and the keeper.

\section{Conclusion}\label{sec:conclusion}

Control of an inbox can be one recovery factor while the chain holds only a
salted commitment: a browser-made DKIM proof over an email whose Subject
blinds the recovery statement, one vote in Perch's delayed recovery.
Soundness, replay resistance and one-factor safety are proven over a model,
privacy in an ideal model; the parsing is tested. Open: the enclave, phones,
real provider mail, people, and an audit before mainnet.


\bibliographystyle{ACM-Reference-Format}
\bibliography{references}

\appendix
\clearpage
\section{Claim ledger}\label{app:claims}

Every claim in this paper and its slides, generated from
\nolinkurl{research/email-recovery/claims.yaml}, in ledger order. A claim is
\claimtag{proven} by a Lean theorem, \claimtag{tested} by a named test,
\claimtag{measured} by a recorded measurement, \claimtag{assumed} as a stated
assumption, or \claimtag{unproven}. Paths are relative to the repository
root.

{\footnotesize\ClaimList}


\end{document}
